\Unit{What this conditional core makes precise}{rat:purpose}{%
FRUIT separates the measured input, a feedback model, a residual-processing
path and a model-rejoin path. A model can alter what enters residual processing
and contribute to the final product. The scientific meaning of that product
therefore depends on the exact path composition, its information sources and
what was fixed or learned along the way.

This revision makes those dependencies explicit without choosing a numerical
procedure. Program-order authority permits authorship; conditional-composition
authority permits deductions under stated premises. Numerical-method authority
is a separate decision. A conditional identity can be exact while every
numerical realization remains unavailable.

The sole current review-candidate normative core is SCI-FRUIT-NORMATIVE-CORE v0.1/r0.4, bound by
the identical source-inventory digest on this cover and the conformance cover.
This rationale explains selected consequences. It does not add requirements,
methods or evidence rules. The complete assumptions, requirements, predictions
and product/state definitions remain in the core.

The frozen Stage A packet and eleven admitted payloads remain unchanged.
The r0.3 fixed-operator repair and three-document structure are accepted.
The r0.4 directive restores only the admitted analysis/gridding-coefficient
ownership and route-use boundary. Internal modal application coordinates remain
separate from those downstream handoff facts. The accepted response
specialization and other scientific rules are unchanged. The r0.4 core is the
current review candidate, without hidden inherited obligations, owner approval,
freeze, Registry registration or activation. Role names are core declarations only.

A later numerical method must resolve the actual parent route, model,
projection, residual processing, rejoin, support, learning, recurrence,
stopping and selector. Upstream producer permissions remain necessary.
Neither this document nor its engineering companion fills those choices.
\ReadCore{Authority and typed objects; REQ-001, REQ-023--024.}
}
\Unit{Three model roles and the applied operand}{rat:models}{%
A candidate model is proposed. An accepted model is selected for feedback. The
applied model is the exact object used by removal and rejoin. Their spaces
may differ, and their relationship is a scientific transformation rather than
an implied equality.

Normalization, clipping, sign conversion, support restriction, projection or
unit conversion can change the applied object. An accepted-model uncertainty
or error therefore cannot be placed directly on applied-model axes without
that transformation and its dependence. Selection and transformation variation
also require separately addressable uncertainty roles.

The conditional composition is imported directly from the normative source:
\EqComposition
Here $\Pi_k(\ma)$ is projected removal, $F_k(\rho_k)$ is the complete
processed residual and $B_k(\ma)$ is rejoin. The method defines every operation
location, bypass and application count. There is no assumed inverse, identical
projection/rejoin, equal input/output space or universal bypass.

Removal and rejoin each require exact compatible quantity, unit/scale, frame,
domain, sampling, support, additive origin/reference or gauge, null-space,
calibration and response conventions. Affine or quotient representations need
well-defined arithmetic in the declared reference; selecting arbitrary
representatives can change the result. Equal array dimensions cannot establish
this compatibility.

A feedback model is not thereby truth, a source detection, posterior or a
Pointing/OOF result. A rejoined unavailable measured mode remains model/prior
supported. The scientific question is which path supplied that content, not
whether a finite number appears in the final array.
\ReadCore{Typed spaces and Composition; REQ-002--005, REQ-015.}
}
\Unit{A hierarchy of exact equations}{rat:equations}{%
The full composition is the starting point. Only a fixed linear complete
residual procedure permits distributive expansion:
\EqMap
This remains map-valued if projection or rejoin is nonlinear. Only complete
fixed linearity of all three maps permits the linear applied-model operator:
\EqLinear
A fixed structural basis alone is insufficient. Centering, operator-defining resolution rules,
support and all consequential state belong to the premise. Affine offsets
remain unless complete linearity is actually established.

The hierarchy prevents a convenient linear expression from silently taking
over a nonlinear or adaptive procedure. The limiting cases are conditional:
zero applied feedback reduces to ordinary residual processing only for
zero-preserving removal/rejoin; identity residual processing cancels removal
and rejoin only when those contributions and spaces agree. If $D_k=0$ under
its complete premises, the fixed result is independent of the applied model.
None of these deductions chooses a numerical operator or proves recovered sky.

For an applied-model difference under the same fixed operators/input, the
result difference is $D_k\delta m_k^{\ap}$. The operator is model-path transfer; the result difference is the distinct
model-induced output shift. Model mismatch and bias have their own identities;
a deterministic difference is not covariance. Calling its expected shift bias needs one exact
estimand, reference/truth model, error random variable, probability law and
conditioning. Without those premises, bias is unavailable even when the shift
is well defined. Accepted error must first be transformed into applied space.
\ReadCore{Composition and model-error transfer; ASM-001--004, REQ-006 and REQ-015.}
}
\Unit{What a response query holds fixed}{rat:response}{%
RF-01 perturbs iteration input; RF-02 perturbs the applied feedback model.
Both freeze complete operator-defining state and the complete induced numerical
maps, including support, reference/gauge/null-space, branch and the other
operand. Internal coordinates that evaluate those maps may change with the
perturbed operand. A changed modal coordinate is not by itself a changed map.

For differentiable fixed maps on compatible tangent spaces:
\EqDerivative
Bind baseline, units/basis, support, codomain, derivative convention and fixed
conditioning. Complete fixed linearity reduces these to $A_k,D_k$; a fixed
affine operation still retains its additive output term. A failed derivative
or domain/comparison failure remains unavailable or a separately typed
finite-difference/transition query, never silently RF-03.

RF-03 reruns prescribed operations that can change defining maps or consequential
state: construction/acceptance, loading subspace, metric, support/mask,
centering/reference, rank, grouping, threshold, selection or branch. Prior
iteration state and upstream procedure/state remain fixed. The distinction
is causal: reevaluating an internal coordinate under a fixed map does not
rerun its defining-state resolution. An unresolved particular map remains
unavailable; calling an internal coordinate an application coefficient
does not change its role.

RF-04 propagates a fixed iteration extent. RF-05 includes stopping and terminal
selection. RF-06 reruns the included upstream chain. RF-07 fixes upstream
producer state/generation and composes the complete FRUIT-only parent-domain
query through the applicable local/recursive/terminal components. Equal values
cannot make these query scopes interchangeable.
\ReadCore{Operator roles; both response sections; ASM-005; REQ-011--013.}
}
\Unit{Why fixed-subspace PTC remains one map}{rat:ptc}{%
PTC owns its estimator and exact fixed-state application map; FRUIT owns only
the response role assigned to a declared use of that operation. The exact
admitted group/time equations yield the following shared core relation:
\EqPTC
Here every symbol is PTC-local. Frozen loading, metric, normal matrix/inverse,
centering, support, detector binding, rank/tolerance, mask, orientation, branch
and other defining facts fix $P$. The modal coordinate
$\widehat m^{\rm app}(x)$ changes algebraically with $x$ while $P$ stays fixed.
This is fixed-map evaluation, not relearning or a new operator generation.
The full FRUIT response still composes its actual remove/process/rejoin paths.

Fixed centering of uncentered calibrated input gives:
\EqPTCAffine
The output is affine; its fixed term $-\lambda_gP$ is retained. Its derivative
is $HP$. A separately named frozen numerical-component subtraction instead
holds $\widehat U$ fixed in $Y'-\widehat U$ and has identity derivative $H$.
Those are distinct operations even if one result happens to agree.

The perturbation must remain in the source's exact finite/rank/support domain.
Failure produces affected unavailability or a separately typed transition.
Rerunning a defining-state rule is a procedure response. Neither outcome
changes the PTC estimator, chooses rank/support or authorizes a numerical
FRUIT method, route or fidelity claim.
\ReadCore{Exact admitted PTC specialization; PRED-006; REQ-012--013.}
}
\Unit{Support and information sources remain visible}{rat:support}{%
Support is not one universal mask. Measured-parent support, the three model
stages' definition support, removal support, residual-processing support,
rejoin support and final-output support describe different scientific roles.
Learning influence records transitive dependence; response-query support
records the exact perturbation and result domain. All eight remain separately
addressable even when some happen to coincide.

The method determines whether compatible residual-and-model, residual-only or
model-only composition is permitted. Missing information is not a numerical
zero. Undefined unequal-support behavior is unavailable, with affected scope
and cause. No implicit fill, extrapolation, renormalization, reference offset
or gauge alignment is supplied by the composition notation.

Occurrence class and mode origin answer different questions. A location can
have both paths available while a particular absolute/additive mode is
supported only by model or prior. Rejoin then adds model-sourced content; it
does not turn that mode into an observationally measured quantity. A reference
choice similarly supplies coordinates, not an observation with zero uncertainty.

Iteration does not create exposure or independent measurements. Reusing a
parent, shared learned model, calibration, response or reference retains
dependence. Iteration count therefore supplies no square-root uncertainty
scaling. Observation-local, coadd-level and joint learning also cannot be
interchanged merely because output shapes match.

These distinctions preserve upstream ownership. In particular, the inherited
JINC base product does not gain a coadd, covariance, significance or response
product because the FRUIT core names a conditional obligation for one.
Unavailable producer authority stays unavailable for its dependent claim.
\ReadCore{Ownership; Path support and information origin; REQ-007--009, REQ-022.}
}
\Unit{Law-relative uncertainty and NOI member meaning}{rat:uncertainty}{%
A stochastic target names its exact law, random variables, fixed objects,
resolution/selection condition, omitted variation and scientific domain/reference.
Pre-outcome means before that law's random event. Physical-parent uncertainty
keeps parent/model values random. Deliberately fixing the realized model defines
a narrower target with learning omitted. A conditional NOI-assignment law can
instead fix realized parent/model/operator/support and exact successful design
resolution before drawing its member. These targets are not equivalent.

Under complete fixed linearity and finite compatible joint moments, the
law-qualified parent/applied covariance has all four terms:
\EqCovariance
Both cross terms remain unless exact law/conditioning justifies omission.
They vanish with model variance only inside the explicitly fixed-model target.
FRUIT branch/completion/convergence/terminal-conditioned targets need separate
method and estimand authority; a generic conditioning symbol supplies none.

With fixed applied model/state, role- and iteration-qualified members obey
\EqNOI
The NOI zero is an expectation under its exact design law, not a demand that
an individual member have zero spatial/sample mean. No finite-ensemble mean
subtraction or $B-1$ correction follows; source/residual structure may remain.
Do not add deterministic rejoin to each zero-target-law member. A full-affine
method needs its own exact center and estimator authority.

The interim owner directive supplies this conditional distinction. Actual
external NOI law/product/estimator binding remains unavailable until supplied;
the notation does not admit a numerical route. Candidate, selection, accepted,
transformation, applied and cross-covariance roles remain independently
addressable, alongside core-declared external/operator/support/selection/empirical
roles. Bias and output shift are distinct products outside that family.
Fixed-state, dependent-successor and relearned NOI remain different methods
and ensembles, with omissions and immutable ancestry explicit.
\ReadCore{Law-relative covariance; NOI and core-declared uncertainty roles; REQ-014--016.}
}
\Unit{Iteration, continuation and terminal selection}{rat:state}{%
The core identifies transition roles without choosing recurrence. The method
must specify whether an iteration consumes the original parent, a previous
result, a reconstruction or another exact product. Every route retains
immutable original measured-parent ancestry. Candidate, accepted, applied,
cumulative model, diagnostic difference and result retain their own meanings.

Unchanged defining-state rules can generate new operator-defining coefficients,
learned states and model values. Those changes require new realized generations,
not a new method each time. Internal fixed-map coordinate evaluation alone
creates no new operator generation.
Changing the learning, support, recurrence, stopping or selection rule changes
method identity. Ordered operations and exact bypass/application counts are
part of that distinction.

After completed absolute zero-based iteration $N$, exact continuation starts
at $N+1$ from $S_{N+1}$. It requires all future-consequential scientific,
stochastic and control state under the same compatible future controls and
numerical profile. A saved result array does not prove sufficiency. Restart
creates a new lineage; retry retains exact attempt and state accounting.

Completion, resource exhaustion, convergence, stopping and terminal selection
are separate claims. A resource limit does not prove convergence, and a
converged sequence does not establish scientific adequacy. No last, latest,
lowest-residual or visually favorable iteration is selected by default.

A selector compares completed candidates in common quantity, units, frame,
domain, gauge, response, support and uncertainty/metric meaning. If conversions
are needed, they and the comparison belong to its method identity. The record
retains every candidate, conversion, inclusion/exclusion, tie and unavailable
comparison. Selection creates an immutable relation to a completed bundle;
it never rewrites the selected iteration.
\ReadCore{Operation state; Completion and terminal comparison; REQ-010--011, REQ-017--020.}
}
\Unit{Local completion and the present claim ceiling}{rat:completion}{%
Each request generation is immutably not-requested or requested; a later
request creates a new generation. For a requested generation, the successful
producer order distinguishes effective action, resolved scope, invocation,
realized path values/outcomes, complete candidate, the FRUIT-owned decision,
completed iteration and publication. Availability and producer outcome are
separate branch/role records at the stage reached; early stop fabricates no
later success. The result is a child payload: scientific consumption requires
its exact bundle or immutable bundle reference.

A complete bundle supplies all universal and method-required roles, or the
precise typed statuses allowed for narrower claims. A method cannot waive core
roles or hard dependencies. Numerical result, continuation, response,
uncertainty and terminal selection each retain the inputs and premises needed
for their own claim. An unavailable optional companion does not erase an
otherwise permitted narrower local result.

Terminal-use references are open-world: record only requested/linked decisions
or an explicit no-request statement. For SCI-VAL use retain exact request,
applicability, eligibility and realization with profile/source binding.
Not-requested creates no eligibility proposition. Missing profile/scope can
produce a realized accounting record with applicability\_unknown and
decision\_unavailable; incomplete/failed/not-produced artifacts are separate
realization conditions. These child decisions do not erase base completeness
or create a FRUIT iteration. Only an explicitly requested narrower product
role may require a child. SCI-VAL evaluates use policy, not FRUIT execution,
model choice or terminal selection.

The engineering specification explains how a future exact candidate would
produce requirement-level evidence. It supplies no numerical conformance now.
No numerical method, route, response/covariance fidelity, recovered measured
mode, coverage, convergence adequacy, observational validation, performance,
readiness, downstream fitness or production authorization is claimed.

The next decision is Grant Wilson's separate review of this revised conditional
core and possible conditional freeze. Numerical development, qualification
and policy recommendations remain separately authorized work with inherited
independent-pointing evidence obligations. The present documents complete
conditional authorship and document verification only.
\ReadCore{Products and lifecycle; Universal completion and assumptions; REQ-018, REQ-021--024.}
}
